\section{ClusterFormer: Token Clustering Attention with Differentiable K-Means and Local Sliding Window}

\textbf{Authors:}
\begin{itemize}
\item {{AUTHOR\_1\_NAME}}  -- {{AUTHOR\_1\_EMAIL}}
\end{itemize}
\begin{itemize}
\item {{AUTHOR\_2\_NAME}}  -- {{AUTHOR\_2\_EMAIL}}
\end{itemize}

\textbf{Affiliations:}
\begin{itemize}
\item   {{UNIVERSITY\_1}}, {{DEPARTMENT\_1}}
\end{itemize}
\begin{itemize}
\item   {{UNIVERSITY\_2}}, {{DEPARTMENT\_2}}
\end{itemize}

\textbf{Corresponding author:} {{CORRESPONDING\_AUTHOR}}
\textbf{Preprint date:} {{DATE}}
\textbf{Version:} v{{VERSION}}
\textbf{Code:} https://github.com/{{ORG}}/SVD\_Project/tree/main/framework/attentions/cluster\_attention.py
\textbf{Model weights:} https://huggingface.co/{{HF\_ORG}}
\textbf{License:} {{LICENSE}}



\subsection{Abstract}

Standard self-attention has quadratic complexity in sequence length, 
$\mathcal{O}(L^2 d)$, which dominates the cost of transformer inference 
beyond a few hundred tokens. Existing linear attention methods achieve 
$\mathcal{O}(L d)$ complexity but sacrifice representational capacity by 
either restricting attention to local windows (Longformer) or approximating 
softmax via random features (Performer). We introduce \textbf{ClusterFormer}, a 
\textbf{differentiable K-means clustering attention} that groups queries into 
$K$ semantic cluster centroids, evaluates \textbf{centroid-to-key cross-attention}, 
and returns to individual queries via a soft assignment. This yields 
$\mathcal{O}(K L d)$ complexity while preserving \textbf{global context}, 
unlike local-only methods.

ClusterFormer combines four components:

1. \textbf{Differentiable K-means} over query representations with padding-aware 
   initialization and 3 Lloyd iterations.
2. \textbf{Centroid-to-key global cross-attention} producing $K$ context vectors 
   at cost $\mathcal{O}(K L d)$.
3. \textbf{Learnable soft assignment} with a temperature parameter to redistribute 
   centroid contexts back to individual queries.
4. \textbf{Local sliding window residual} blended with the global branch via a 
   learnable gate, capturing nearby syntactic structure that clustering may 
   miss.

Base Q/K/V/O projections are frozen and copied from the pretrained model; 
LoRA rank-16 adapters and learnable temperature/blend scalars provide the 
only trainable capacity. On DistilBERT, RoBERTa-base, and GPT-2 with SST-2 
as the evaluation task, ClusterFormer maintains accuracy within 
\textbf{{{D\_DELTA}}\%--{{G\_DELTA}}\%} of the baseline after {{RECOVERY\_STEPS}} 
recovery steps while reducing attention FLOPs by \textbf{{{FLOP\_REDUCTION}}$\times$} at 
sequence length $L = 128$ with $K = {{NUM_CLUSTERS}}$ clusters.

\textbf{Keywords:} linear attention, differentiable clustering, token routing, 
sliding window, efficient transformers



\subsection{Lay Summary}

Standard transformer attention compares every token to every other token, 
which becomes slow when the sequence gets long. ClusterFormer takes a 
different approach: it groups similar tokens into a small number of 
"clusters" (like sorting people by hometown at a conference) and then 
lets each cluster attend to all tokens at once. This is much faster, and 
it still keeps global context -- every cluster sees the whole sequence. 
To capture fine-grained local structure, we also add a small sliding window 
that attends to nearby tokens. The two signals are mixed together with a 
learnable gate. On three popular language models, ClusterFormer matches 
the original models' accuracy on sentiment analysis while using a fraction 
of the attention computation.



\subsection{Table of Contents}

1. [Introduction](01\_intro.md)
2. [Related Work](02\_related.md)
3. [Background](03\_background.md)
4. [Method](04\_method.md)
5. [Architecture Integration](05\_architecture.md)
6. [Experimental Setup](06\_setup.md)
7. [Results](07\_results.md)
8. [Analysis](08\_analysis.md)
9. [Discussion](09\_discussion.md)
10. [Limitations](10\_limitations.md)
11. [Conclusion](11\_conclusion.md)
12. [Reproducibility](12\_reproducibility.md)
13. [References](13\_references.md)

\textbf{Appendices:}
\begin{itemize}
\item [A. Mathematical Derivations](A\_math.md)
\end{itemize}
\begin{itemize}
\item [B. Implementation Details](B\_implementation.md)
\end{itemize}
\begin{itemize}
\item [C. Full Experimental Data](C\_data.md)
\end{itemize}
\begin{itemize}
\item [D. Negative Results](D\_negative\_results.md)
\end{itemize}
\begin{itemize}
\item [E. Ethical Considerations](E\_ethics.md)
\end{itemize}